\documentclass[runningheads]{llncs}

% ICECCS 2026 short paper (max 11 pages incl. references), Springer LNCS format.
% https://formal-analysis.com/iceccs/2026/

\usepackage[T1]{fontenc}
\usepackage[utf8]{inputenc}
\usepackage{graphicx}
\usepackage{booktabs}
\usepackage{array}
\usepackage{xcolor}
\usepackage{listings}
\usepackage{url}
\usepackage{amsmath}
\usepackage[hidelinks]{hyperref}

\graphicspath{{figures/}}

\begin{document}

\title{An AI-Aware Security Gateway for Retrieval-Augmented Chatbots: A Cross-Model Evaluation}
\titlerunning{An AI-Aware Security Gateway for RAG Chatbots}

% TODO: confirm affiliations, ORCIDs, and per-author emails before submission.
\author{%
Daniel Martins Ferreira\inst{1}\and
Rui Machado\inst{1}\and
Henrique Teixeira\inst{1}\and
Sebasti\~ao Teixeira\inst{1}\and
Osvaldo Pacheco\inst{1}}
\authorrunning{D.\,M.\,Ferreira et al.}

\institute{University of Aveiro, Aveiro, Portugal\\
\email{danielmartinsferreira@ua.pt}}

\maketitle

\begin{abstract}
Retrieval-augmented generation (RAG) lets chatbots answer from an
organisation's own documents, but it widens the attack surface: untrusted
retrieved content and user input reach the model and its tools directly,
exposing prompt injection, data leakage, and tool abuse. A common defence is
to rely on the underlying model's alignment, yet that defence varies by model
and is hard to audit. This paper studies an alternative: an AI-aware security
gateway that enforces structural controls (input filtering, retrieval
authorisation, context quarantine, output redaction, and scoped tool use)
independently of the model. We harden an open-source RAG security tool into a
deployable gateway, integrate it with a retrieval-augmented chatbot over a real
tourism-services knowledge base, and evaluate it against four OWASP LLM Top-10
attack classes. For each of several open-source models we measure a
guardrails-off baseline and compare it with the gateway enabled, separating
model alignment from model-independent structural controls.
% TODO: add one-sentence headline result once experiments are run.

\keywords{RAG security \and prompt injection \and LLM guardrails \and OWASP LLM Top-10 \and AI gateway \and open-source LLMs.}
\end{abstract}

\section{Introduction}
\label{sec:intro}
% Motivation: organisations deploy RAG chatbots over private corpora; the
% security model is immature. State the gap: alignment-based defences are
% model-dependent and unauditable. State the contribution.
%
% Contributions (short-paper scope):
%  C1. An AI-aware gateway hardened from an open-source research tool into a
%      reusable, drop-in component ("more in-the-box").
%  C2. A reference RAG chatbot deployed over the ATT tourism-services corpus.
%  C3. A cross-model evaluation against four OWASP LLM Top-10 attack classes,
%      with a guardrails-off baseline per model.
\textbf{TODO:} write introduction (motivation, gap, contributions, paper map).

\section{Background and Related Work}
\label{sec:background}
% RAG (Lewis et al.); prompt injection direct/indirect (Perez \& Ribeiro;
% Greshake et al.); OWASP Top 10 for LLM Applications; existing guardrail /
% gateway approaches and their limits. Position our model-independent,
% structural-controls angle.
\textbf{TODO:} background on RAG security and the OWASP LLM Top-10; related guardrail work.

\section{The AI-Aware Security Gateway}
\label{sec:gateway}
% Describe the gateway as a Policy Enforcement Point on the RAG pipeline:
% Input Filter -> Retrieval (authorise as user) -> Quarantine -> LLM ->
% Output Filter -> Tools (scoped + validated). Explain the hardening from the
% research prototype into a deployable, configurable component, and the
% threat model (the four attack classes S1-S4).
\textbf{TODO:} architecture figure + description of the five structural controls and the engineering work to make the tool deployable.

\section{Deployment over a Tourism-Services Knowledge Base}
\label{sec:deployment}
% The reference RAG chatbot built on the ATT microservices platform
% (knowledge-base service as corpus). Entitlement model (clearance levels),
% retrieval, and where the gateway sits.
\textbf{TODO:} chatbot design, corpus, user entitlement model, integration points.

\section{Evaluation}
\label{sec:evaluation}
% Method: scenarios x {vulnerable, secure} x models x runs. Models: list the
% open-source models compared (e.g. via Ollama). Baseline = guardrails off per
% model. Metric: per-scenario attack-success score; separate model alignment
% (vulnerable-mode refusals) from structural controls (secure mode).
\subsection{Setup}
\textbf{TODO:} models under test, attack scenarios, metrics, runs, hardware.

\subsection{Results}
% Table: per-model attack-success rate, vulnerable vs secure, per scenario.
\textbf{TODO:} results table + key findings.

\section{Discussion and Limitations}
\label{sec:discussion}
\textbf{TODO:} what the structural controls do and do not cover; threats to validity; generalisation beyond the tourism corpus.

\section{Conclusion and Future Work}
\label{sec:conclusion}
\textbf{TODO:} conclusion and next steps.

\subsubsection*{Acknowledgements.}
% TODO: funding / project acknowledgements (ATT project, etc.).

\bibliographystyle{splncs04}
\bibliography{refs}

\end{document}
